% Source-derived chapter 08: Counting points over finite fields
% Original source: no-enriques-surface-over-the-integers
\chapter{Counting points over finite fields}
\label{no-enriques-surface-over-the-integers-chapter-08}
\source{no-enriques-surface-over-the-integers}

\begin{remark}[Source section]
\label{no-enriques-surface-over-the-integers-chapter-08-source}
\source{no-enriques-surface-over-the-integers}
\source{no-enriques-surface-over-the-integers}
This chapter reproduces the corresponding section of the retrieved
LaTeX source, including its definitions, statements, examples, and
proofs. It has not yet been translated into Lean.
\notready
\end{remark}

% The source section title was: Counting points over finite fields
\mylabel{Counting points}

In this section we count the number of rational points  on Enriques surfaces  with constant local system over finite fields.
The observation has nothing in particular to do with Enriques surfaces, so we first work in a general setting.
Let $k=\FF_q$ be a finite ground field, for some prime power $q=p^\nu$, and   $X$ be a smooth proper scheme of dimension $n\geq 0$,
with $h^0(\O_X)=1$.

Set $N_s=\Card X(\FF_{q^s})$ for  the integers $s\geq 1$, such that $N_1$ is the number of rational points.
The \emph{Hasse--Weil zeta function} is 
\begin{equation}
\label{zeta function}
\source{no-enriques-surface-over-the-integers}
Z(X,t) = \prod_a \frac{1}{1-t^{\deg\kappa(a)}} =\exp\left(\sum_{s=1}^\infty \frac{N_s }{s}t^s\right),
\end{equation}
where the product runs over all closed points $a\in X$. This formal power series actually belongs to
the subring $\QQ(t)\subset\QQ[[t]]$, by Dwork's Rationality Theorem \cite{Dwork 1960}.

Choose some algebraic closure $\bk=k^\alg$, write $\bX=X\otimes \bk$, and fix a prime $\ell\neq p$. 
For each $i\geq 0$ and each integer $j$ we can form the  $\ell$-adic cohomology groups
\begin{equation}
\label{l-adic cohomology groups}
\source{no-enriques-surface-over-the-integers}
H^i(\bX,\QQ_\ell(j))= \left( \invlim_r H^i(\bX,\mu^{\otimes j}_{\ell^r})\right)\otimes_{\ZZ_\ell}\QQ_\ell.
\end{equation}
Here $j\in\ZZ$ refers to the \emph{Tate twist} in the coefficient sheaves. Note that $H^i(\bX,\QQ_\ell(j))$
arises from $H^i(\bX,\QQ_\ell)$ by tensoring with the one-dimensional  vector space
$\QQ_\ell(j)=(\invlim_r\mu_{l^r}(\bk))\otimes_{\ZZ_\ell}\QQ_\ell$, which comprises only information on the action of $\Gal(\bk/k)$
on the $l$-primary roots of unity in $\bk$.
In particular, the $\QQ_\ell$-dimensions $b_i\geq 0$ of the $\ell$-adic cohomology 
$H^i(\bX,\QQ_\ell(j))$ do not depend on the Tate twist, and are called  the \emph{Betti numbers} of $X$.

The   power $F^\nu_X:X\ra X$ of the absolute Frobenius is a $k$-morphism, and induces by base-change a 
$\bk$-morphism $\Phi=F^\nu_X\otimes \id_\bk$ on $\bX$. Let us write $f_{i,j}=\Phi^*$ for the induced 
$\QQ_\ell$-linear endomorphism on the    cohomology group $H^i(\bX,\QQ_\ell(j))$. Of particular interest are the $f_i=f_{i,0}$:
The zeta function takes the form 
$$
Z(X,t)=\prod \det(1-tf_i)^{(-1)^{i+1}},
$$
by the  Grothendieck--Lefschetz Trace Formula \cite{FL}.
Here the  polynomial factors can be seen as
``reciprocal'' characteristic polynomials 
$\det(1-f_it)=t^{b_i}\det(t^{-1}-f_i)$ of the endomorphisms $f_i$. 
The characteristic polynomials of $f_{i,j}$ have coefficients from $\QQ$, and the eigenvalues
$\alpha\in\QQ^\alg$ are algebraic integers, 
of  length $|\alpha|=q^{i/2-j}$ in all complex embeddings, according to the Riemann Hypothesis (\cite{Deligne 1980}, Corollary 3.3.9,
see   \cite{Deligne 1974}, Theorem 1.6 for the projective case).

The following observations explain the  effect of the Tate twist:
First of all, we have a canonical identification $H^i(\bX,\QQ_\ell(j))=H^i(\bX,\QQ_\ell)\otimes \QQ_\ell(j)$.
Moreover, the  composition $(F_X^\nu\otimes\id_\bk)\circ (\id_X\otimes F_\bk^\nu)$ is a power of the absolute Frobenius of $\bX$, which
acts trivially on \'etale cohomology   with respect to any coefficient sheaf.
Finally, the Frobenius power $F_\bk^\nu$ acts by multiplication with $q=p^\nu$ on $\QQ_\ell(1)$. Combining these three observations,
we see that $f_{i,j}=f_i\otimes q^{-j}$, which reveals that the eigenvalues of $f_i$ are obtained from the eigenvalues of $f_{i,j}$
by multiplication with $q^j$.

The Tate twists become   relevant in connection with algebraic cycles:
Let $\CH^j(\bX)$ be the \emph{Chow group} of $j$-codimensional cycles modulo linear equivalence, which comes with 
a \emph{cycle class map}  $\CH^j(\bX)\ra H^{2j}(\bX,\QQ_\ell(j))$. 
We likewise can form Chow groups and  $\ell$-adic cohomology on the scheme $X$, before base-changing to the algebraic closure,
and have a commutative diagram 
\begin{equation}
\label{cycle class maps}
\source{no-enriques-surface-over-the-integers}
\begin{CD}
\source{no-enriques-surface-over-the-integers}
\notready
\CH^j(X)	@>>> 	H^{2j}(X,\QQ_\ell(j))\\
@VVV			@VVV\\
\CH^j(\bX)	@>>> 	H^{2j}(\bX,\QQ_\ell(j))
\end{CD}
\end{equation}
of cycle class maps. For more details on cycle class maps see \cite{Deligne 1977}.

\begin{proposition}
\source{no-enriques-surface-over-the-integers}
\notready
\mylabel{generalization wedderburn}
Suppose for all integers $j\geq 0$ we have $H^{2j+1}(\bX,\QQ_\ell)=0$, and that the 
composite map $\CH^j(X)\otimes\QQ_\ell \ra H^{2j}(\bX,\QQ_\ell(j))$ is surjective.
Then the zeta function is
\begin{equation}
\label{simple zeta function}
\source{no-enriques-surface-over-the-integers}
Z(X,t) = \prod_{j=0}^{n}\left(\frac{1}{1-q^jt}\right)^{b_{2j}}.
\end{equation}
In particular, the number of rational points is $\Card X(\FF_q)= \sum_{j=0}^{n} b_{2j}q^j$, and  the set of rational points is non-empty.
\end{proposition}

\proof
Suppose  that the endomorphism $f_{2j}$ is the homothety with eigenvalue $q^j$. Then 
$$
t^{b_{2j}}\det(t^{-1}-f_{2j}) = t^{b_{2j}}(t^{-1}-q^j)^{b_{2j}} = (1- q^jt)^{b_{2j}},
$$
and the assertion on the zeta function follows. Expanding \eqref{zeta function} and \eqref{simple zeta function} 
as formal power series in $t$ and comparing linear terms, we get the statement on the number of rational points.
Since $b_0=1$ the set  $X(\FF_q)$ must be non-empty.

It remains to verify the assertion on the endomorphism $f_{2j}$. 
Set $r=b_{2j}$ and choose prime cycles $Z_1,\ldots,Z_r\subset X$ whose cycle classes in $H^{2j}(\bX,\QQ_\ell(j))$ form
a vector space basis. Obviously, the action of $\id_X\otimes F_\bk$ on the base changes $\bar{Z}_1,\ldots,\bar{Z}_r\subset\bX$ is trivial,
so  the same holds for the action on $H^{2j}(\bX,\QQ_\ell(j))$. The latter action is   inverse to the action of
$\Phi=F_X^\nu\otimes\id_\bk$, as observed above.  In other words, $f_{2j,j}$ is the identity map.
We also observed above that $f_{2j,j}=f_{2j}\otimes q^{-j}$, thus $f_{2j}$ is the homothety with eigenvalue $q^j$.
\qed

\medskip
The conditions obviously hold if $b_{2j+1}=0$ and $b_{2j}=1$ for all $j\geq 0$. One may see the above result
as a generalization of  Wedderburn's Theorem, which states that every Brauer--Severi variety $X$ contains a rational point, 
and is thus isomorphic to   projective  space $\PP^n$. 
Let us now specialize to surfaces. Recall that $\rho=\rank \Num(\bX)$ denotes the Picard number.

\begin{corollary}
\source{no-enriques-surface-over-the-integers}
\notready
Suppose that $X$ is a surface with $b_1=0$ and $b_2=\rho$.
If the local system $\Num_{X/k}$ is constant, we have $\Card X(\FF_p)= 1+b_2q+q^2$. 
\end{corollary}

\proof
We have $b_0=1$, and with Poincar\'e Duality get $b_4=1$ and $b_1=0$. The cycle class map $\CH^1(X)\ra H^2(\bX,\QQ_\ell(1))$
factors over $\Num(X)$, and  by the assumptions  the inclusion $\Num(X)\otimes\QQ_\ell\subset H^2(\bX,\QQ_\ell(1))$ is an equality.
Thus we may apply the  Proposition and  the formula on  $\Card X(\FF_p)$ follows.
\qed

\begin{corollary}
\source{no-enriques-surface-over-the-integers}
\notready
\mylabel{points on enriques}
Suppose that $X$ is either an Enriques surface, or a  geometrically rational   surface endowed
with a relatively minimal genus-one fibration $X\ra\PP^1$. If the local system $\Num_{X/k}$ is constant,
we have  $\Card X(\FF_q) = 1+10q+q^2$.
\end{corollary}

\proof
In both cases, we have $b_1=0$. 
In light of the previous corollary, we only have to show
that $\bX$ has Betti number $b_2=10$. For Enriques surfaces, this holds by definition.
In the other case, the induced  fibration over $\bk$ remains a relatively minimal genus one-fibration.
We have $K_{\bX}^2=0$ by \cite{Cossec; Dolgachev 1989}, Proposition 5.6.1.
Hence there is a morphism $\bX\ra\PP^2_\bk$ that can be factored as a sequence of nine blowing ups,
so we also have $b_2=10$.   
\qed

\medskip
We shall be particularly concerned with the case $q=p=2$. The above formula then gives $\Card X(\FF_2)=25$.


%===========================================================
